\section{What the original specifies, and what it does not}
\label{sec:original}

This section sets out the index of \citet{lagasio2024} as published, then delimits where
the published specification stops. The exercise is documentary: it uses none of our
pipeline and asserts nothing about what the original's code did. Where the article is
silent we record the silence; we do not substitute what standard practice would have been.
There is one exception, and because it is the pivot on which a referee would press we make
it in the open rather than in a footnote: the \texttt{norm} setting is inferable from the
convention that only non-default library arguments are reported, we make that inference and
rely on it in \S\ref{sec:measurement}, and we accordingly do not count it among the gaps
that make this index irreproducible.
The article is open access under CC BY~4.0, so contested points are quoted rather than
paraphrased. Page numbers are the article's own.

\subsection{The published construction}

The index has two components. The sentiment component $\SEN_i$ is the score returned by the
TextBlob library for report $i$ (p.~4), which the index definition calls the ``sentiment
polarity score''. The substance component $\SUS_i$ is derived from scikit-learn's
\texttt{TfidfVectorizer}, version 0.24.2 (p.~3), evaluated against an author-constructed
list of ESG keywords. The index is the difference of the two components
after score normalisation, and the labelling rule is definitional: $\ESGSI > 0$ is
``Potential ESGwashing'', $\ESGSI \le 0$ is ``Likely Genuine'' (pp.~4--5). No calibration
against an external benchmark is reported.

The vectoriser settings are disclosed in two sentences (p.~4):

\begin{quote}
``We use Scikit-learn's TfidfVectorizer with settings chosen to balance comprehensiveness
and computational efficiency, in line with best practices in NLP research. These settings
include ignoring terms that appear in more than 95~\% of documents, only considering terms
that appear in at least 2 documents, limiting the vocabulary to the top 10,000 terms, and
considering both unigrams and bigrams.''
\end{quote}

In library terms, \texttt{max\_df=0.95}, \texttt{min\_df=2}, \texttt{max\_features=10000},
\texttt{ngram\_range=(1,2)}. These four are the only vectoriser settings the article
discloses, and they are described in prose only: the article prints no parameter names, no
code listing and no equation for $\SUS_i$. Two further sentences in the same paragraph bear
on the vectoriser and add no settings. The first establishes that IDF weighting is in use:
the method evaluates ``the relative importance of each term within the documents against the
entire dataset'' (p.~4). The second states that the vectoriser ``is used with adjustable
settings to accommodate the unique features of different sustainability reports'' (p.~4);
no setting beyond the four above is given.

The step from per-term values to a document-level score is where the specification stops.
Two sentences bear on it. From the original's \S3.2.1 (p.~4): ``Each document is then
vectorized with the TF-IDF model to calculate a score for each term based on its frequency
relative to other documents.'' From its \S3.2.3 (p.~4): ``It combines sentiment scores with TF-IDF scores of
sustainability terms to objectively evaluate how extensively sustainability topics are
discussed in the reports.'' The abstract adds only ``integrating sentiment analysis with
the frequency of sustainability terms to calculate the ESGSI'' (p.~1). In the index formula,
$\SUS_i$ is glossed as ``sustainability score for company i'' --- a name, not a
construction.

\subsection{What is disclosed, and what is not}

Table~\ref{tab:original-disclosure} separates the two.

\begin{table}[ht]
\centering
\small
\caption{Elements of the original specification, by disclosure status.}
\label{tab:original-disclosure}
\begin{tabular}{p{0.46\textwidth}p{0.44\textwidth}}
\toprule
Element & Status in \citet{lagasio2024} \\
\midrule
Vectoriser implementation and version & Disclosed: scikit-learn \texttt{TfidfVectorizer} 0.24.2 (pp.~3--4) \\
\texttt{max\_df}, \texttt{min\_df}, \texttt{max\_features}, \texttt{ngram\_range} & Disclosed in prose (p.~4) \\
IDF weighting in use & Disclosed (p.~4) \\
Sentiment instrument & Disclosed: TextBlob polarity (p.~4) \\
Keyword-list provenance & Disclosed: GRI (2021) and SASB (2018) standards (p.~4) \\
Labelling threshold & Disclosed, definitional (pp.~4--5) \\
\addlinespace
\texttt{norm} (vector normalisation) & Not stated; library default reasonably inferable \\
\texttt{sublinear\_tf}, \texttt{smooth\_idf} & Not stated; library defaults reasonably inferable \\
Aggregation from per-term TF-IDF to $\SUS_i$ & Not stated; not a library setting, so no default supplies it \\
Whether TF-IDF is computed over the 10,000-term vocabulary and then subset to the keywords, or the vectoriser is restricted to the keywords & Not stated \\
Keyword-list size and contents & Not published; three example terms only \\
Custom stopword additions to the NLTK defaults & Not published \\
Whether raw or preprocessed text reached TextBlob & Not stated \\
Score normalisation (min--max or $z$-score) & Contradictory; see \S\ref{sec:artefact} \\
Share of firms flagged ``Potential ESGwashing'' & Never reported \\
Cohen's Kappa, 5-fold cross-validation, TF-IDF sensitivity analysis & Announced (p.~5), never reported \\
Data availability statement & ``No'' (p.~10) \\
\bottomrule
\end{tabular}
\end{table}

Four entries warrant expansion. First, \emph{vector} normalisation. The \texttt{norm}
parameter is never named; the strings \emph{L2} and \emph{Euclidean} do not occur in the
article at all; and no rescaling of document vectors is described under any other wording.
The article does use the word \emph{normalization}, but of two other operations: lowercasing
and character stripping during preprocessing (p.~3), and the rescaling of the finished
component scores across the sample (p.~4), which is a different operation at a different
stage (\S\ref{sec:artefact}). Table~\ref{tab:original-disclosure} therefore records
\texttt{norm} as not stated, which is a fact about the text; it does not record it as a
reproducibility gap, and we concede the point openly. The four disclosed settings are
exactly four non-default \texttt{TfidfVectorizer} arguments, so a reader following the
ordinary convention that only departures from library defaults are reported would infer
that the remaining defaults --- \texttt{norm='l2'} among them --- were in force. That
inference is reasonable, we rely on it ourselves, and \S\ref{sec:measurement} draws it
explicitly and labels it as an inference. It remains an inference rather than a statement
of the article, so we do not say that the original L2-normalised its document vectors and
we do not say that it did not; but a silence that a standard convention closes is not where
the reproducibility charge of \S\ref{sec:original-consequence} rests.

Second, the vocabulary. Its provenance is given --- the indicator lists ``are based on
established ESG reporting frameworks such as Global Reporting Initiative (GRI) Standards
\ldots and Sustainable Accounting Standards Board (SASB) Standards'' (p.~4) --- but its
size is never stated and the list is never published. Three example terms appear
(``Diversity'', ``Carbon Emissions'', ``Equal Employment Opportunities''); there is no
appendix, table or supplementary file.

Third, the announced validation (p.~5):

\begin{quote}
``For the manual data collection process, we employ two independent raters and calculate
Cohen's Kappa to assess agreement \ldots We use 5-fold cross-validation to evaluate the
robustness of our LDA and sentiment analysis models \ldots We also conduct sensitivity
analyses by varying key parameters (e.g., number of LDA topics, TF-IDF settings) to test
the robustness of our results.''
\end{quote}

Neither a Kappa value nor a cross-validation figure nor a sensitivity result appears
anywhere in the article: outside the reference list, the passage above is the only
occurrence of \emph{Kappa}, of \emph{cross-validation} and of \emph{sensitivity} in the
text. The announced sensitivity analysis names TF-IDF settings explicitly, so a reported
version of it would bear directly on the vectoriser silences recorded above
(\S\ref{sec:measurement}).

Fourth, the data availability statement, which consists of the heading ``Data
availability'' and the single word ``No'' (p.~10).

\subsection{One further disagreement, resolved elsewhere}

Within its \S3.3 the article prints a formula that is a difference of $z$-scores ---
reproduced verbatim in \S\ref{sec:index} --- and states in prose, in the same subsection,
that ``both sentiment and sustainability scores are normalized to a uniform scale using
min-max normalization'' (p.~4). These are different estimators and the article never
reconciles them. We record the disagreement here and adjudicate it in
\S\ref{sec:artefact}, on evidence of a different kind; nothing in the present section turns
on which reading is correct.

\subsection{Consequence for what follows}
\label{sec:original-consequence}

A reader who set out to compute this index would have to supply, from outside the article,
at minimum: the aggregation rule from per-term values to $\SUS_i$, the keyword list and its
size, the score-normalisation rule, and the text passed to the sentiment instrument. Each
is a free parameter, and \S\ref{sec:robustness} shows that admissible settings differ in the
numbers they produce. The index as published is therefore not reproducible from its own
description. The obstacle is not access --- the full text was available to us --- but the
completeness of the description.

The first of those four is the strongest instance and carries the charge on its own. The
step from a vector of per-term TF-IDF values to a single $\SUS_i$ is a modelling decision,
not a library argument: sum, mean, maximum and weighted variants are all available, the
article names none of them, and no reporting convention supplies a default because the
choice is not a default of anything. This is the respect in which it differs from
\texttt{norm}, where the convention that only non-default arguments are reported does close
the gap and we have said so above. The second is of the same kind: a keyword list that is
neither published nor sized cannot be reconstructed by any convention, and it is what makes
strict replication impossible at the lexicon stage (\S\ref{sec:lexicon}). The third is not a
silence at all but a contradiction --- the prose of the original's \S3.3 and the formula
printed in the same subsection specify different estimators, and no convention adjudicates
between two statements the article itself makes. None of these three is answerable from
outside the article, and the charge stands on them.

This governs the design of what follows. Because the original specification does not
determine a unique implementation, we do not implement one: we implement several and report
results across them. That is why \S\ref{sec:index} carries three substance specifications
rather than one, and why \S\ref{sec:robustness} reports how the index and its labels move
between them.
